\section{Algoritma Metaheuristik untuk Optimasi Kombinatorial}

Optimasi kombinatorial merupakan proses mencari solusi terbaik dari ruang solusi diskret yang sangat besar dan umumnya bersifat \textit{NP-hard}, sehingga penyelesaian eksak menjadi tidak praktis untuk instansi berskala besar.
Metaheuristik menawarkan pendekatan pencarian tingkat tinggi untuk memperoleh solusi mendekati optimal dengan biaya komputasi yang lebih rendah, meskipun tanpa jaminan optimalitas formal.
Perkembangan bidang ini sangat pesat; \textcite{rajwar_exhaustive_2023} mencatat lebih dari 500 algoritma metaheuristik telah dikembangkan, dengan lebih dari 350 di antaranya muncul dalam satu dekade terakhir.

Berdasarkan sumber inspirasinya, metaheuristik dapat dikelompokkan menjadi algoritma evolusioner, seperti \textit{Genetic Algorithm} dan \textit{Differential Evolution}; \textit{swarm intelligence}, seperti \textit{Particle Swarm Optimization} dan \textit{Ant Colony Optimization}; algoritma berbasis fenomena fisika atau kimia, seperti \textit{Simulated Annealing} dan \textit{Gravitational Search Algorithm}; serta kategori lain yang terinspirasi perilaku manusia, permainan, atau konsep matematis \parencite{rajwar_exhaustive_2023}.
Namun, klasifikasi metaheuristik belum sepenuhnya konsisten antar-penulis dan perlu terus diperbarui seiring munculnya metode baru \parencite{jaksic_comprehensive_2023}.

Dari sisi mekanisme pencarian, metaheuristik juga dapat dibedakan menjadi \textit{trajectory-based} dan \textit{population-based}.
Algoritma \textit{trajectory-based}, seperti \textit{Simulated Annealing} dan \textit{Tabu Search}, menggunakan satu agen yang bergerak melalui ruang solusi, sedangkan algoritma \textit{population-based}, seperti \textit{Genetic Algorithm, Particle Swarm Optimization}, dan \textit{Ant Colony Optimization}, menggunakan banyak agen secara simultan untuk memperluas eksplorasi ruang solusi \parencite{rajwar_exhaustive_2023}.

ACO dan UBK yang digunakan sebagai \textit{baseline} dalam penelitian ini termasuk metaheuristik \textit{population-based}.
ACO meniru perilaku koloni semut dalam menemukan jalur, sedangkan UBK terinspirasi strategi berburu \textit{black-winged kite}.
Keduanya menggunakan sekumpulan kandidat solusi untuk mengeksplorasi ruang kebijakan ABAC secara simultan.

Meskipun meningkatkan kemampuan eksplorasi, pendekatan \textit{population-based} membutuhkan penyimpanan banyak kandidat solusi secara eksplisit selama proses optimasi.
Akibatnya, kebutuhan memori meningkat seiring ukuran populasi.
Karakteristik ini menjadi kendala pada perangkat \textit{edge-IoT} dengan kapasitas memori terbatas dan menjadi dasar kebutuhan terhadap pendekatan optimasi yang lebih \textit{compact}.